% Section 9 — Limitations and outlook.
% Reuses the ECCV "Limitations" subsection but expands with software /
% deployment limitations relevant to GMD readers.

We highlight the limitations of the current release so that prospective
users and contributors can decide whether CRAN-PM is appropriate for
their use case.

\paragraph{Spatial coverage.}
The pre-trained Europe model covers roughly $30^\circ$N--$72^\circ$N,
$25^\circ$W--$45^\circ$E. Out-of-distribution evaluation on USA, Canada
and India (Section~\ref{sec:eval}, Fig.~\ref{fig:robinson}) shows useful
zero-shot transfer for some regions, but operational use outside Europe
should be preceded by region-specific fine-tuning, for which we provide
the \texttt{cranpm.training.finetune} entry point.

\paragraph{Pollutants.}
v1.0 forecasts only PM\textsubscript{2.5}. Extending to PM\textsubscript{10},
NO\textsubscript{2} and O\textsubscript{3} is straightforward
architecturally — only the output channel needs to change — but requires
the corresponding reanalysis training targets.

\paragraph{Lead time.}
The model is trained for 1--4 day lead times. Extending to seasonal
forecasting would require a different training setup (lower temporal
resolution, larger context window).

\paragraph{Operational input dependencies.}
Inference at time $t$ requires: ERA5 (latency $\sim$5 days from
Copernicus), CAMS analysis (low latency), and the GHAP daily PM2.5
analysis at $t-1$ (longer latency, $\sim$2 weeks). Real-time deployment
therefore requires either a near-real-time GHAP product or a
self-bootstrapping autoregressive variant. We are exploring both as
future work.

\paragraph{Uncertainty.}
The current release outputs a deterministic forecast. A simple
Monte-Carlo dropout estimate is provided as
\texttt{cranpm.inference.uncertainty}, but a properly calibrated ensemble
(multi-seed or perturbed-input) is left for v0.2.

\paragraph{Compute requirements.}
End-to-end training requires roughly 64 GPU-days on MI250X / A100. This
is moderate by deep-learning standards but precludes laptop-only training.
Inference is much cheaper: a full European map at 1\,km is produced in
\textasciitilde 1.8\,s on a single GPU.

\paragraph{Reproducibility caveats.}
The MIOpen workaround in the patch-embedding (Section~\ref{sec:software})
ensures bit-identical results between CUDA and ROCm at \texttt{bf16}.
Mixed-precision training nonetheless retains a small run-to-run
non-determinism arising from non-deterministic atomic reductions in
attention kernels. We document the residual variance in
\texttt{tests/test\_reproducibility.py}.
